% This file is part of the StegX project.
% Copyright (C) 2018  StegX Team
% 
% StegX is free software: you can redistribute it and/or modify
% it under the terms of the GNU General Public License as published by
% the Free Software Foundation, either version 3 of the License, or
% (at your option) any later version.
% 
% This program is distributed in the hope that it will be useful,
% but WITHOUT ANY WARRANTY; without even the implied warranty of
% MERCHANTABILITY or FITNESS FOR A PARTICULAR PURPOSE.  See the
% GNU General Public License for more details.
% 
% You should have received a copy of the GNU General Public License
% along with this program.  If not, see <https://www.gnu.org/licenses/>.

\documentclass[11pt]{article}
%\documentclass{book}
\usepackage[utf8]{inputenc}
\usepackage[T1]{fontenc}
\usepackage[french]{babel}
\usepackage[top=1.8cm, bottom=1.8cm, left=1.8cm, right=1.8cm]{geometry}
\usepackage[linktocpage,colorlinks=false]{hyperref}
\usepackage{graphicx}
\usepackage{epsfig}
\usepackage{amssymb}
\usepackage{amsmath}
\usepackage{array}
\usepackage{subfig}
\usepackage{multicol}
\usepackage{caption}
\usepackage{listings}
\usepackage{algorithm}
\usepackage{algorithmic}
\usepackage{array,multirow,makecell}
\hypersetup{
    colorlinks=true,
    breaklinks=true,
    urlcolor=red,
}
\parskip=5pt

\title{\huge{\textbf Ajouts du compte rendu version du 1{er} juin}}
\author{AYOUB Pierre, BASKEVITCH Claire, BESSAC Tristan, \\
CAUMES Clément, DELAUNAY Damien, DOUDOUH Yassin}
\date{Mercredi 1 Juin 2018}

\begin{document}

\maketitle
\vspace{20em}
\begin{center}\includegraphics{pictures/Application.png}\end{center}
\newpage

\section{Introduction}

L'équipe StegX a profité de la semaine jusqu'au 1er juin pour améliorer 
l'application. De ce fait, sachant que le client ne proposait aucun format 
à s'occuper, les développeurs ont voulu ajouter le format MP3. 
Pourtant, le MP3 est un format très compliqué étant compressé. 

Pour la version du 25 mai, Pierre Ayoub et Damien Delaunay avaient réussi à 
étudier les versions du format (MPEG 1 Layer III, MPEG 2 Layer III) et de 
format de métadonnée (ID3 version 1 et ID3 version 2). 
Et ils ne leur restaient plus qu'à implémenter un ou plusieurs algorithmes 
pour le format MP3. 
De ce fait, MP3 propose désormais les algorithmes EOF et LSB. 

Les études approfondies sur le MP3 ont mené à modifier certains choix dans 
les spécifications en raison de la grande difficulté de ce format. 

\section{Changements des spécifications}

\subsection{Structure MP3}

\begin{lstlisting}[language=c]
struct mp3 {
    long int fr_frst_adr; 
    long int fr_nb;       
    long int eof;      
};
\end{lstlisting}

Cette structure comporte 3 champs : 
\begin{itemize}
\item \textit{fr\_frst\_adr} est l'adresse du header de la première frame 
MPEG 1/2 Layer III. 
\item \textit{fr\_nb} est le nombre de frame MPEG 1/2 Layer III.
\item \textit{eof} est l'adresse de la fin du fichier officiel (sans signature et données cachées).
\end{itemize}

Cette structure n'a pas pu être déduite lors des spécifications puisque 
l'étude détaillée du format était trop bas niveau à cette période du projet. 
En effet, la stéganographie sur le format MP3 est très expérimentale puisque 
ce format est très complexe. 

\subsection{Ajouts de fonctions pour éviter la répétition de code}

Pour réaliser correctement le remplissage de la structure afin d'éviter la
repétition de code lors des algorithmes proposés par le format MP3, 
nous avons créé plusieurs fonctions internes au format : 

\begin{lstlisting}[language=c]
int mp3_mpeg_hdr_test(uint32_t hdr);
int mp3_id3v2_hdr_test(uint32_t hdr);
\end{lstlisting}

Ces fonctions testent si le header hdr est un header MPEG 1/2 Layer III
ou un header ID3v2.
Elles renvoient 0 si le header est incorrect et 1 si le header est valide. 
\newline

\begin{lstlisting}[language=c]
int mp3_mpeg_fr_seek(uint32_t hdr, FILE * f);
long int mp3_mpeg_fr_find_first(FILE * f);
int mp3_mpeg_hdr_get_version(const uint32_t hdr);
int mp3_mpeg_hdr_is_padding(const uint32_t hdr);
int mp3_mpeg_hdr_get_samprate(const uint32_t hdr);
int mp3_mpeg_hdr_get_bitrate(const uint32_t hdr);
int mp3_mpeg_hdr_get_size(const uint32_t hdr);
int mp3_mpeg_fr_write(uint32_t hdr, FILE * src, FILE * dst);
\end{lstlisting}

Ces fonctions concernent les tests internes au format MP3 et font les tests 
en rapport aux blocs MPEG. 
\newline

\begin{lstlisting}[language=c]
int mp3_id3v1_hdr_test(uint32_t hdr);
int mp3_id3v1_tag_seek(FILE * f);
int mp3_id3v1_tag_write(FILE * src, FILE * dst);
\end{lstlisting}

Ces fonctions concernent les tests internes au format MP3 et font les tests 
en rapport aux blocs IDv1. 
\newline
\begin{lstlisting}[language=c]
uint32_t mp3_id3v2_size_unsyncsafe(const uint32_t s);
int mp3_id3v2_tag_seek(FILE * f);
\end{lstlisting}

Ces fonctions concernent les tests internes au format MP3 et font les tests 
en rapport aux blocs IDv2. 
\newline \newline
Le listing détaillé de ces fonctions est 
présent dans les fichiers \textit{mp3.h/.c}.

\subsection{Ajouts dûs à la résolution d'un bug}

Lors de la préparation de la démonstration, nous voulions montrer l'aspect 
multiplateforme de l'application. 
Ainsi, nous avons remarqué que, lorsque l'on fait une insertion sur un 
système d'exploitation, l'extraction doit se faire sur le même système
d'exploitation (pas forcément sur la même machine). 
Ce bug est dû au fait que la suite pseudo aléatoire générée sur Linux n'est 
pas la même que sur Windows. 
De ce fait, il a fallu créer nos propres fonctions de randomisation pour que 
la suite pseudo-aléatoire soit la même pour n'importe quelle système d'exploitation : 

\begin{lstlisting}[language=c]
void stegx_srand(unsigned int seed);
\end{lstlisting}

Cette fonction va initialiser la graine de la suite pseudo aléatoire du 
système. \textit{seed} correspond au nombre où la graine sera initialisée.
\newline 

\begin{lstlisting}[language=c]
int stegx_rand();
\end{lstlisting}

Cette fonction va renvoyer un entier pseudo-aléatoirement. 

\section{Bilan technique du produit pour la version du 1er juin}

En ce qui concerne les algorithmes proposés par StegX en fonction des 
formats pris en charge, voici le bilan de ce que propose l'application 
après la version du 1er juin: 
\newline

\begin{tabular}{|l|c|c|c|c|c|}
  \hline
  \multirow{2}*{\textbf{Format pris en charge par l'application}} & \multicolumn{5}{c|}{\textbf{Algorithmes proposés}} \\
   \cline{2-6}
    & \textbf{LSB} & \textbf{EOF} & \textbf{Metadata} 
    &\textbf{EOC} & \textbf{Junk Chunk} \\
  \hline
  \textbf{BMP} (Clément Caumes \& Yassin Doudouh) & \textbf{\checkmark} & \textbf{\checkmark} & \textbf{\checkmark} &  & \\
  \hline      
  \textbf{PNG} (Clément Caumes \& Yassin Doudouh) &  & \textbf{\checkmark} & \textbf{\checkmark} & & \\
  \hline
  \textbf{WAV} (Pierre Ayoub \& Damien Delaunay) & \textbf{\checkmark} & \textbf{\checkmark} & & & \\
  \hline 
  \textbf{MP3} (Pierre Ayoub \& Damien Delaunay) & \color{red}{\textbf{\checkmark}} & \color{red}{\textbf{\checkmark}} & & & \\
  \hline 
  \textbf{AVI} (Claire Baskevitch \& Tristan Bessac) & & & & & \textbf{\checkmark}\\
  \hline
  \textbf{FLV} (Claire Baskevitch \& Tristan Bessac) & & \textbf{\checkmark} & & \textbf{\checkmark} & \\
  \hline
\end{tabular}
\vspace{0.5cm}

On peut remarquer que la ligne pour le format MP3 est remplie. StegX propose 
maintenant de la stéganographie sur les formats MP3. 

\end{document}
